\section{Proof of the central theorem}\label{sec:proof}
\begin{lemma}[Actual action mass]\label{lem:mass}
For a regenerative task and any admissible history procedure, set
\[
 x_a=\E\sum_{t=1}^{\tau}\one_{\{A_t=a\}},\qquad
 \mu_a(D)=\Pp(A_\tau=a,X_\tau\in D).
\]
Then $\mu_a\le x_a\nu_a$, $\sum_a\mu_a(K_a)=1$, and $\sum_a c_ax_a$ is the actual expected cost. Under \eqref{eq:envelopes}, $\mu_{\theta,a}\le\overline b_a x_{\theta,a}\nu_a$.
\end{lemma}
\begin{proof}
The event $\{\tau\ge t\}$ is measurable before the $t$th call. Condition also on its selected action, and use \eqref{eq:regen} to obtain
\[
 \E\sum_{t=1}^{\tau}\one_{\{A_t=a,X_t\in D\}}
       =\nu_a(D)\E\sum_{t=1}^{\tau}\one_{\{A_t=a\}}.
\]
Tonelli applies to these nonnegative sums. The terminal count is bounded by the acquired count, which proves domination. Positive minimum cost and finite cost imply $\E\tau<\infty$, hence terminal mass one. The same calculation with the upper envelope proves the last assertion. No law after a rejection has been conditioned back to its original value.
\end{proof}

\begin{lemma}[Exact causal realization]\label{lem:realization}
Every feasible point of \eqref{eq:profile} has a one-live-state realization with exact action counts and exact action-indexed terminal submeasures.
\end{lemma}
\begin{proof}
Let $T=\sum_a x_a\ge\sum_a\int h_a=1$. At the single live label choose action $a$ with probability $x_a/T$. If $x_a>0$, after response $z$ stop with probability $h_a(z)/x_a$, and route to a nearest center with a fixed Borel tie rule. Otherwise continue in the same live label. An action of count zero is never chosen.

At every surviving call the stopping probability is exactly $1/T$. Its action-indexed accepted submeasure is $h_a\nu_a/T$. Thus $\tau$ is geometric of mean $T$, the expected number of action-$a$ calls is $x_a$, and the terminal submeasure is $h_a\nu_a$. The expected cost is exactly $\sum_a c_ax_a$. No random value persists after a rejection; fresh action and routing coins are used. Only one live and $m$ terminal labels are present. This also implements feasible points with positive counts but zero accepted mass for some action.
\end{proof}

\begin{lemma}[Existence and the mass--cost dual]\label{lem:dual}
The primal and dual statements in Theorem~\ref{thm:central}\textup{(i)--(ii)} hold.
\end{lemma}
\begin{proof}
The counts lie in the compact box $0\le x_a\le N/c_a$. The constraints on $h_a$ are weak-star closed in the corresponding bounded subsets of $L^\infty(\nu_a)$. For example, $h_a\le x_a$ is equivalent to all nonnegative $L^1$ integral inequalities and is preserved when the scalar $x_a$ converges. The centers lie in compact $H^m$. Their squared distance functions converge uniformly as the centers converge, so the objective is jointly continuous for these topologies. A minimum follows. Feasibility is equivalent to $N\ge\min c_a$: necessity uses $\sum x_a\ge1$, and sufficiency takes one cheapest action with $x=h=1$.

For fixed $C$, attach multipliers $z\in\R$ to unit terminal mass and $\lambda\ge0$ to the cost bound. Minimizing the Lagrangian over $0\le h_a\le x_a$ gives
\[
 z-\lambda N+\sum_a x_a\left\{\lambda c_a-\int(z-D_C)_+\,d\nu_a\right\}.
\]
Its infimum over $x_a\ge0$ is $z-\lambda N$ precisely under the inequalities in \eqref{eq:dual}. This proves weak duality.

For completeness, strong duality is a finite-dimensional separation argument, not a minimax exchange with the centers. Consider the convex cone of triples (terminal mass, cost, distortion) obtained from $x_a\ge0$ and $0\le h_a\le x_a$, with nonnegative slack in the last two coordinates. On every bounded-cost slice its generating variables are weak-star compact, so the cone is closed locally where the cost is bounded. At mass one and cost $N>c_*:=\min c_a$, the lower distortion boundary has a supporting affine plane, since the projection onto (mass,cost) contains the open cone $\{u>0,v>c_*u\}$. Its cost coefficient is nonpositive, because increasing a permitted cost cannot increase the value. Normalize the distortion coefficient to one; this is possible at an interior projected point where the finite convex value function has a subgradient. Writing the plane as distortion $\ge zu-\lambda v$ gives $\lambda\ge0$. Testing it on all rays from one action gives exactly the dual inequalities. Equality at $(1,N,V_C)$ proves attainment.

At $N=c_*$ take $N_k\downarrow c_*$. Compactness of primal minimizers shows $V_C(N_k)\uparrow V_C(c_*)$. A dual optimizer at $N_k$ evaluated at $c_*$ has value $V_C(N_k)+\lambda_k(N_k-c_*)\ge V_C(N_k)$. Weak duality and passage to the limit prove the endpoint supremum formula. This argument does not require bounded endpoint multipliers.
\end{proof}

\begin{proof}[Proof of Theorem~\ref{thm:central}\textup{(i)}]
The protected conditional projection identity \eqref{eq:pythagoras} and Lemma~\ref{lem:mass} bound every checkpoint with $m$ output values below by $B+\Psi_m$. An autonomous observer has at most $M-1$ output values, since a live state is necessary. Thus it is bounded below by $B+\Psi_{M-1}$. Lemmas~\ref{lem:realization} and \ref{lem:dual} give attaining programs with exactly the allowed architecture, proving both equalities.

We prove the asserted two-action choice directly. Start at a primal optimizer and discard an action with zero accepted mass by setting its count to zero; this cannot increase cost. Let $w_a=\int h_a\,d\nu_a>0$ on the remaining actions. Freeze the conditional terminal law $h_a\nu_a/w_a$, the per-unit terminal cost $v_a=c_ax_a/w_a$, and its distortion $e_a=\int D_Ch_a\,d\nu_a/w_a$. The original weights satisfy $\sum w_a=1$ and $\sum v_aw_a\le N$. Minimize $\sum e_aw_a$ over this finite polytope. An extreme minimizer has at most two positive weights: with three positive weights there is a nonzero direction preserving both total mass and cost, and both sufficiently small signs remain feasible, contradicting extremality. If the cost is slack the same conclusion, or the stronger one-action conclusion, follows. Rescaling each frozen count and density by its new weight preserves $h_a\le x_a$. The new objective cannot exceed the global optimum, so it is again optimal and Lemma~\ref{lem:realization} realizes it. This elementary support reduction is consistent with, but does not require, general constrained-policy reduction theorems~\cite{DPRchat}.
\end{proof}

\begin{proof}[Proof of Theorem~\ref{thm:central}\textup{(iii)}]
Fix a model and any legal observer. Lemma~\ref{lem:mass} gives terminal densities bounded by $\overline b_a x_{\theta,a}$. Put $x'_a=\overline b_a x_{\theta,a}$. Their cost with coefficients $c_a/\overline b_a$ is at most $N$. This proves the two lower bounds for each model and hence for the common worst-model optimum.

For the upper, choose a minimizing point for $\Psi_m(N;c/\underline b)$. Put $T=\sum_a x_a/\underline b_a\ge1$. At one live state choose $a$ with probability $x_a/(\underline b_aT)$ and accept response $z$ with probability $h_a(z)/x_a$. At every surviving history its stopping probability is at least $1/T$. Therefore $\E_\theta\tau\le T$. Its mean cost per surviving call is the fixed number $\sum_a c_ax_a/(\underline b_aT)$, giving total mean cost at most $N$. Its conditional distortion contribution on a call is at most
\[
 \frac1T\sum_a\frac{\overline b_a}{\underline b_a}
                 \int D_Ch_a\,d\nu_a
 \le\frac\kappa T\Psi_m(N;c/\underline b).
\]
Sum over the actual survival probabilities to prove the upper. The argument is uniform over the unknown fixed parameter and all conditional rows, including those after rejection. It does not claim that the accepted law equals its regenerative reference law.
\end{proof}

\begin{proof}[Proof of Theorem~\ref{thm:central}\textup{(iv)}]
Use an optimal one-live-state program for the unprimed experiment. Its action probabilities and report routing induce a joint one-call law on the response, the continuation symbol, and terminal labels. Write $s\ge1/N$ for the one-call stopping mass and $A$ for the subprobability measure of accepted response--label pairs. The risk in excess of $B$ is $\int D\,dA/s$, where $D$ is the squared loss for that routing. The primed one-call law has accepted submeasure $A'$ and stopping mass $s'$, with $|s-s'|\le\eta$. Since $D\one_{\{\text{stop}\}}\in[0,L]$, its numerator changes by at most $L\eta$.

If $s'<s$, redirect a fraction $(s-s')/(1-s')$ of continuing outcomes to one already existing terminal label. This is a fresh routing coin, costs no new label, and adds accepted mass $d=s-s'\le\eta$. If $s'\ge s$, make no change. The repaired stop mass is $s''=\max(s,s')$, so its mean is $1/s''\le1/s\le N$. Its excess risk is at most
\[
 \frac{\int D\,dA+L\eta+L(s-s')_+}{s''}
 \le \frac{\int D\,dA}{s}+2LN\eta.
\]
A stop mass $s'=1$ requires no repair. The construction is valid at $N=1$ as well. Interchanging the two known experiments proves \eqref{eq:linear}; the diameter bound gives the truncation at $L$. The repair uses the known primed one-call mass, not an unknown parameter or an unbounded-cost consequence of transcript proximity.
\end{proof}

\begin{proof}[Proof of Corollary~\ref{cor:regular}]
For a union of $m$ balls of radius $r$ about arbitrary centers, \eqref{eq:regular} gives $\nu_a$-mass at most $C m r^{s_a}$; if a center is outside $K_a$, move it to a point of intersection and double the radius. Lemma~\ref{lem:mass} and unit cost show that the actual terminal mass in this union is at most
\[
 C mN\max_a r^{s_a}.
\]
Choose $r=c\min_a(Nm)^{-1/s_a}$, with $c$ small enough for all exponents and for $r\le r_0/2$. At least half the terminal mass lies outside the balls. This proves the lower bound, with $m=M-1$ for autonomous observers.

For the upper use only a minimizing action. We recall the required global argument. A maximal $r$-separated subset of an $s$-regular support gives a global cover with at most $C r^{-s}$ points, by the lower mass bound on disjoint balls. The upper mass bound, conversely, produces $m$ separated points at distance $c m^{-1/s}$. For $t$ above a fixed constant, their disjoint radius $C(tm)^{-1/s}$ balls have total reference mass at least $1/t$. Accept a normalized restriction to that union; its density is at most $t$ and distortion at most $C(tm)^{-2/s}$. For bounded $t\ge1$ use the whole-law global cover instead. These statements retain the actual mass and prove $\Psi_m^{\nu_a}(t)\asymp(tm)^{-2/s_a}$. Apply the envelope upper with $t=\underline b_a N$ when this is at least one. When $1\le N<1/\underline b_a$, stop after one call and use the global cover. Constants are uniform over the finite action family. No trajectorywise contraction or asymptotic dimension theorem is inserted into this construction.
\end{proof}
